% Section: Sumsets in the integers, survey Subsections 3.3--3.4 (continued from integers_A).
% Label tag: IA:
%
% Local conventions: B+B and B+mB are unrestricted sums (b=b' allowed), as in the survey;
%   B \oplus mB = {b+mb' : b,b' in B, b \neq b'} denotes the restricted version.

\section{Ordered polynomial sumsets}\label{IA:sec:ordered}

In this section we answer two unnumbered questions from~\cite[\S3.3]{kra2023problems}. In
\cref{IA:sec:u33a} we determine exactly which ordered polynomial sums
$\{f(b_1)+g(b_2):b_1<b_2\}$ appear, up to a shift, in every set of positive density: those with
$\deg f\le1$ and $\deg g\le1$. In \cref{IA:sec:u33e} we observe that the natural generalisations of
Question~3.18 of~\cite{kra2023problems} have negative answers. Both answers rely on work of
Ackelsberg~\cite{ackelsberg2024counterexamples} and Kousek~\cite{kousek2025asymmetric}; our
contribution is small and we say precisely what it is.

\subsection{Which ordered polynomial sums are unavoidable}\label{IA:sec:u33a}
For functions $f,g$ and $B\subset\N$ write
\[
f(B)\oplus_<g(B)=\{f(b_1)+g(b_2):\ b_1,b_2\in B,\ b_1<b_2\}.
\]
The survey considers these ``ordered sums'' for polynomials $f,g$ ``over $\N$''
\cite[(5)]{kra2023problems}, notes that they are partition regular by Ramsey's theorem, and asks:

\begin{question}{\S3.3(a)}\label{IA:q:u33a}
Let $f$ and $g$ be polynomials over $\N$. Can instances of $f(B)\oplus_<g(B)$ be found, up to a
shift, in every set with positive density? That is, does every $A\subset\N$ of positive density
contain $f(B)\oplus_<g(B)+t$ for some infinite $B\subset\N$ and some $t\ge0$?
\end{question}

The survey does not specify the density; its default is upper Banach density. Special cases are
Theorem~2.2 of~\cite{kra2023problems} ($f=g=x$, positive, in the restricted form), Conjecture~3.10
($f,g$ linear), Question~3.13 ($f=x^2$, $g=x$) and Example~3.14 ($g=x^2$). We interpret ``polynomials
over $\N$'' in the widest sensible way and let
\[
\mathcal P=\{f\in\Q[x]:\ f(\N)\subset\N_0\};
\]
this contains all polynomials with coefficients in $\N_0$.

\begin{theorem}\label{IA:thm:u33a}
Let $f,g\in\mathcal P$.
\begin{enumerate}[label=\textup{(\alph*)}]
\item\label{IA:it:u33a-pos} If $\deg f\le1$ and $\deg g\le1$, then for every $A\subset\N$ with
$\ubd(A)>0$ there are an infinite $B\subset\N$ and $t\ge0$ with $f(B)\oplus_<g(B)+t\subset A$.
\item\label{IA:it:u33a-neg} If $\max(\deg f,\deg g)\ge2$, then for every $\eps>0$ there is a set
$A\subset\N$ with $\ldens(A)>1-\eps$ such that $f(B)\oplus_<g(B)+t\not\subset A$ for every
infinite $B\subset\N$ and every $t\in\Z$. Moreover, one set $A$ works simultaneously for all such
pairs $(f,g)$.
\end{enumerate}
\end{theorem}

Since the hypothesis in~\ref{IA:it:u33a-pos} is the weakest density condition and the conclusion
in~\ref{IA:it:u33a-neg} concerns the strongest one, the classification does not depend on which
notion of density is used in \cref{IA:q:u33a}.

\begin{remark}[Attribution]\label{IA:rem:u33a-credit}
Part~\ref{IA:it:u33a-pos} is a direct consequence of Kousek's theorem~\cite[Theorem~1.2]{kousek2025asymmetric}
(which proved Conjecture~3.10 of~\cite{kra2023problems}); we only add the bookkeeping for constant
terms and zero coefficients. Part~\ref{IA:it:u33a-neg} was known for $\deg g\ge2$ in the reversed or
sparse forms of~\cite[Example~3.14]{kra2023problems}, and for $g(x)=x$ by Ackelsberg's
theorem~\cite[Theorem~1.5]{ackelsberg2024counterexamples}. Our proof of~\ref{IA:it:u33a-neg} is
Ackelsberg's argument with $P(x)+x$ replaced by $f(x)+g(x)$; the only cases not covered explicitly
in the literature that we know of are those with $\deg f\ge2$ and $g(x)=mx+c$, $m\ge2$.
\end{remark}

\begin{remark}[Scope]\label{IA:rem:u33a-scope}
\begin{enumerate}[label=(\roman*)]
\item The argument for~\ref{IA:it:u33a-neg} does not use the order $b_1<b_2$: the same set $A$
also avoids $\{f(b_1)+g(b_2):b_1>b_2\}+t$ and $\{f(b_1)+g(b_2):b_1\neq b_2\}+t$.
\item Integer polynomials that are not $\N_0$-valued on $\N$ are outside \cref{IA:q:u33a}, and the
classification does not extend to them. For instance, for $f(x)=-x$ and $g(x)=x$ the configuration is
the infinite set $\{b_2-b_1:b_1<b_2\}$ and $f+g=0$; since $|R_k(0)|=1$, the proof
of~\ref{IA:it:u33a-neg} below shows that it can be avoided by sets of lower density $>1-\eps$, although
both polynomials are linear.
\item The survey remarks that it is immaterial whether one uses $f(B)\oplus_<g(B)$ or
$\{r+s:r\in f(B),s\in g(B),r<s\}$. This is not accurate in general. If $g$ is non-constant and
$f(b)\ge g(b)$ for all large $b$, the two formulations are equivalent: refining $B$ we may assume
$f(b_1)<g(b_2)$ for all $b_1<b_2$ in $B$, and also that $f$ is non-decreasing and $f\ge g$ on $B$, so
that $f(b_1)<g(b_2)$ forces $b_1<b_2$. If instead $f(b)<g(b)$ for all large $b$, the second set
contains the diagonal elements $f(b)+g(b)$, and the answer can change even for linear $f,g$: for
$f(x)=x$, $g(x)=2x$ the second set contains $\{b_1+2b_2:b_1\le b_2\}$, and
Kousek~\cite[Proposition~5.1]{kousek2025asymmetric} (with $m=1$, $\ell=2$) gives a set of positive
upper density containing no shift $\{b_1+2b_2:b_1\le b_2\}+t$, $t\in\N$, with $B$ infinite (the
case $t=0$ reduces to $t=3$ by replacing $B$ with $B-1$). If $g\equiv c$ and $\deg f\ge2$, the
second set is finite, while $f(B)\oplus_<g(B)=f(B\setminus\{\max B\})+c$ is infinite.
\Cref{IA:thm:u33a} uses the formulation~\cite[(5)]{kra2023problems}.
\item In $\Q$ the situation is different: Ackelsberg~\cite[Theorem~1.4]{ackelsberg2025infinite}
proved that if $P\in\Q[x]$ and $P(x)+x$ are non-constant, then every $A\subset\Q$ of positive upper
Banach density contains $\{b_i,\ P(b_i)+b_j:i<j\}+t$ for some infinite $\{b_n\}\subset\Q$ and some
$t\in\Q$.
\end{enumerate}
\end{remark}

We collect elementary facts about $\mathcal P$.

\begin{lemma}\label{IA:lem:P}
Let $f,g\in\mathcal P$.
\begin{enumerate}[label=\textup{(\roman*)}]
\item $f(\Z)\subset\Z$, and there is $D\in\N$ with $Df\in\Z[x]$. For such $D$ and every $k\in\N$,
$n\equiv n'\pmod{Dk}$ implies $f(n)\equiv f(n')\pmod k$.
\item If $f$ is non-constant, its leading coefficient is positive. Consequently
$\deg(f+g)=\max(\deg f,\deg g)$.
\item If $\deg f\le1$, then $f(x)=mx+c$ with $m\in\N_0$ and $c\in\Z$.
\end{enumerate}
\end{lemma}

\begin{proof}
(i) Write $f(x)=\sum_{i}\Delta^if(1)\binom{x-1}{i}$ (Newton's forward difference formula); the
numbers $\Delta^if(1)$ are integers because $f(\N)\subset\Z$, and $\binom{x-1}i\in\Z$ for $x\in\Z$.
A common denominator $D$ of the coefficients exists since $f\in\Q[x]$. If $n-n'=Dks$, then
$D(f(n)-f(n'))=(n-n')Q$ for some $Q\in\Z$ (as $Df\in\Z[x]$), so $f(n)-f(n')=ksQ$.
(ii) A non-constant polynomial with negative leading coefficient tends to $-\infty$. If
$\deg f=\deg g$, the leading coefficients are both positive and do not cancel.
(iii) $m=f(2)-f(1)\in\Z$, $c=f(1)-m\in\Z$, and $m\ge0$ by (ii).
\end{proof}

For an integer-valued $h\in\Q[x]$ and $k\in\N$ let $R_k(h)=\{h(n)\bmod k:n\in\Z\}\subset\Z/k\Z$.

\begin{lemma}[Ackelsberg]\label{IA:lem:Rk}
Let $h$ be a polynomial in $\Q[x]$ of degree at least $2$ such that $h(\Z)\subset\Z$. Then
$\inf_{k\in\N}|R_k(h)|/k=0$.
\end{lemma}

\begin{proof}
Let $\P_h$ be the set of primes $p$ such that $h(\Z)$ meets every residue class modulo $p$. By
Ackelsberg's Lemma~4.3 in~\cite{ackelsberg2024counterexamples}, $\sum_{p\in\P\setminus\P_h}1/p=\infty$,
so $\prod_{p\in\P\setminus\P_h}(1-1/p)=0$. If $p_1,\dots,p_s\in\P\setminus\P_h$ are distinct and
$k=p_1\cdots p_s$, then by the Chinese remainder theorem every element of $R_k(h)$ reduces modulo
$p_i$ to an element of $R_{p_i}(h)$, so $|R_k(h)|\le\prod_i|R_{p_i}(h)|\le\prod_i(p_i-1)$ and
$|R_k(h)|/k\le\prod_i(1-1/p_i)$. This is the argument of the proof
of~\cite[Theorem~1.5]{ackelsberg2024counterexamples}, applied to $h$ instead of $P(x)+x$.
\end{proof}

\begin{lemma}\label{IA:lem:subadd}
Let $E_1,E_2,\ldots\subset\N$ have natural densities $\delta_1,\delta_2,\dots$. There is
$S\subset\N$ such that $E_j\setminus S$ is finite for every $j$ and $\udens(S)\le\sum_j\delta_j$.
\end{lemma}

\begin{proof}
This is the natural-density version of~\cite[Lemma~3.1]{ackelsberg2024counterexamples}. Choose
$M_1<M_2<\cdots$ such that $|E_i\cap[N]|\le(\delta_i+2^{-J})N$ for all $J$, all $i\le J$ and all
$N\ge M_J$, and let $S=\bigcup_j\big(E_j\cap(M_j,\infty)\big)$. Then $E_j\setminus S\subset[M_j]$.
If $M_J\le N<M_{J+1}$, then $S\cap[N]\subset\bigcup_{j\le J}(E_j\cap[N])$, so
$|S\cap[N]|\le(\sum_j\delta_j+J2^{-J})N$.
\end{proof}

\begin{proof}[Proof of \cref{IA:thm:u33a}\ref{IA:it:u33a-neg}]
The set of pairs $(f,g)\in\mathcal P^2$ with $\max(\deg f,\deg g)\ge2$ is countable; enumerate the
triples $(f,g,t)$ with such $(f,g)$ and $t\in\Z$ as $(f_j,g_j,t_j)_{j\ge1}$. Put $h_j=f_j+g_j$. By
\cref{IA:lem:P}, $h_j$ is integer-valued on $\Z$ and $\deg h_j\ge2$, so by \cref{IA:lem:Rk} there is
$k_j$ with $|R_{k_j}(h_j)|/k_j<\eps2^{-j}$. Let
\[
E_j=\{x\in\N:\ (x-t_j)\bmod k_j\in R_{k_j}(h_j)\},
\]
a union of residue classes of natural density $\delta_j=|R_{k_j}(h_j)|/k_j$. Let $S$ be given by
\cref{IA:lem:subadd} and $A=\N\setminus S$. Then $\ldens(A)=1-\udens(S)\ge1-\sum_j\delta_j>1-\eps$.

Suppose that $f_j(B)\oplus_<g_j(B)+t_j\subset A$ for some $j$ and some infinite $B$. Let $k=k_j$ and
let $D$ be such that $Df_j,Dg_j\in\Z[x]$. By the pigeonhole principle there is $n$ such that
$B'=B\cap(n+Dk\Z)$ is infinite. By \cref{IA:lem:P}(i), $f_j(b_1)+g_j(b_2)\equiv f_j(n)+g_j(n)=h_j(n)
\pmod k$ for all $b_1,b_2\in B'$, so $f_j(B')\oplus_<g_j(B')+t_j\subset A\cap E_j\subset E_j\setminus
S$, which is finite. But $f_j(B')\oplus_<g_j(B')$ is infinite: if $g_j$ is non-constant, fix
$b_1=\min B'$ and let $b_2\to\infty$; if $g_j$ is constant, then $\deg f_j\ge2$ and $f_j(b_1)\to\infty$
as $b_1\to\infty$ in $B'$. This contradiction proves~\ref{IA:it:u33a-neg} for every $(f,g)$ and every
$t\in\Z$ simultaneously.
\end{proof}

\begin{proof}[Proof of \cref{IA:thm:u33a}\ref{IA:it:u33a-pos}]
By \cref{IA:lem:P}(iii), $f(x)=mx+c$ and $g(x)=\ell x+c'$ with $m,\ell\in\N_0$ and $c,c'\in\Z$. Let
$\ubd(A)>0$; in particular $A$ is infinite.

If $m,\ell\ge1$, Kousek's theorem~\cite[Theorem~1.2]{kousek2025asymmetric} gives an infinite $B$ and
$t\in\N$ with $\{mb_1+\ell b_2:b_1,b_2\in B,\,b_1<b_2\}+t\subset A$. For $j\in\N$ let
$B_j=\{b-j:b\in B,\,b>j\}$. Then
\[
f(B_j)\oplus_<g(B_j)+t_j'=\{mb_1+\ell b_2:\ b_1,b_2\in B,\ j<b_1<b_2\}+t\subset A,
\qquad t_j'=t-c-c'+(m+\ell)j,
\]
and $t'_j\ge0$ for $j$ large.

If $m=0<\ell$, then $f(B)\oplus_<g(B)=\{\ell b+c+c':b\in B\setminus\{\min B\}\}$. Choose $r$ such
that $A\cap(r+\ell\Z)$ is infinite, then $r'\equiv r\pmod\ell$ with $r'\ge c+c'$, and put $t=r'-c-c'$
and $B=\{(x-r')/\ell:\ x\in A,\ x\equiv r\ (\mathrm{mod}\ \ell),\ x>r'\}$. Then $\ell b+c+c'+t\in A$ for
all $b\in B$. The case $\ell=0<m$ is the same, since then
$f(B)\oplus_<g(B)=\{mb+c+c':b\in B\}$ (every element of an infinite $B$ has a larger one). If
$m=\ell=0$, the configuration is the single point $c+c'$; take $t=x-c-c'$ for some $x\in A$ with
$x\ge c+c'$.
\end{proof}

\formalization{\Cref{IA:thm:u33a} is formalized in \leanfile{S33a}. Part~\ref{IA:it:u33a-neg}
(\leanname{thm_u33a_neg}) is unconditional; there Ackelsberg's Lemma~4.3, used in the proof of
\cref{IA:lem:Rk}, is proved by a different argument (Hermite's criterion and Dirichlet's
theorem). Part~\ref{IA:it:u33a-pos} (\leanname{thm_u33a_pos}) takes Kousek's
theorem~\cite[Theorem~1.2]{kousek2025asymmetric} as an explicit hypothesis.}

\subsection{Generalisations of Question~3.18}\label{IA:sec:u33e}
The survey asks whether shifts of
$\{b_i^2+b_j:i<j\}$ appear in every set of positive upper Banach density
\cite[Question~3.13]{kra2023problems}, and generalises this by adding the base point $b_j$
\cite[Question~3.15]{kra2023problems} and by replacing $x^2$ with arbitrary polynomials
$p_1,\dots,p_k\in\Z[x]$ with $p_\ell(0)=0$ \cite[Question~3.16]{kra2023problems}. Its
Remark~3.17 records that the answers to Questions~3.15 and~3.16 are negative by Ackelsberg's work;
indeed, both configurations contain that of Question~3.13, which Ackelsberg showed can be
avoided~\cite[Theorem~1.5]{ackelsberg2024counterexamples}. Then it asks:

\begin{question}{3.18}\label{IA:q:318}
If $A\subset\N$ has positive upper Banach density, must there exist an infinite set
$B=\{b_1<b_2<\cdots\}$ and $t\ge0$ such that $\{b_ib_j+b_k:\ i<j<k\}$ is contained in $A-t$?
\end{question}

\begin{question}{\S3.3(e)}\label{IA:q:u33e}
Similar to the way Question~3.13 is generalised in Questions~3.15 and~3.16, one can inquire about
analogous generalisations of Question~3.18.
\end{question}

The faithful analogues are:
\begin{itemize}
\item[(F1)] (base point, as in Question~3.15) must every $A$ with $\ubd(A)>0$ contain
$\{b_k,\ b_ib_j+b_k:\ i<j<k\}+t$ for some infinite $B$ and $t\ge0$?
\item[(F2)] (polynomial family, as in Question~3.16) given $p_1,\dots,p_r\in\Z[x,y]$ with
$p_\ell(0,0)=0$, must every $A$ with $\ubd(A)>0$ contain
$\{p_\ell(b_i,b_j)+b_k:\ 1\le\ell\le r,\ i<j<k\}+t$ for some infinite $B$ and $t\ge0$?
\end{itemize}
As in Remark~3.17 of~\cite{kra2023problems}, (F2) is a question about all such families.

\begin{proposition}\label{IA:prop:u33e}
Both (F1) and (F2) have negative answers. More precisely, for every $\eps>0$ there is $A\subset\N$
with $\ubd(A)>1-\eps$ such that for no $t\in\Z$ and no infinite $B=\{b_1<b_2<\cdots\}\subset\N$ does
$A-t$ contain the configuration in (F1), or the configuration in (F2) for the one-element family
$p_1(x,y)=xy$.
\end{proposition}

\begin{proof}
Ackelsberg~\cite[Theorem~1.7]{ackelsberg2024counterexamples} proved that for every $\eps>0$ there is
$A\subset\N$ with $\ubd(A)>1-\eps$ such that whenever $t\in\Z$ and
$\{b_1b_2+b_3+t:\ b_1,b_2,b_3\in B,\ b_1<b_2<b_3\}\subset A$, the set $B$ is finite; this answers
\cref{IA:q:318} negatively. The configuration in (F1), and the configuration in (F2) for $p_1=xy$,
both contain $\{b_ib_j+b_k:i<j<k\}$, so this set $A$ has the required property.
\end{proof}

\formalization{\Cref{IA:prop:u33e} and the remark following it are formalized in
\leanfile{S33e} (\leanname{prop_u33e}, \leanname{F2_sq}), including proofs of the two quoted
theorems of Ackelsberg in the form used here.}

For (F2) one does not even need Question~3.18: for the one-element family $p_1(x,y)=x^2$, if
$\{b_i^2+b_k:i<j<k\}+t\subset A$, then $C=\{b_1,b_3,b_5,\dots\}$ satisfies
$\{c^2+c':c,c'\in C,\ c<c'\}+t\subset A$ (between $b_{2i-1}$ and $b_{2i'-1}$, $i<i'$, lies $b_{2i}$),
which is the configuration of Question~3.13 of~\cite{kra2023problems}; by Ackelsberg's
theorem~\cite[Theorem~1.5]{ackelsberg2024counterexamples}, recorded in~\cite{kra2023problems},
there are sets of upper Banach density arbitrarily close to $1$ in which no translate contains it.

\begin{remark}\label{IA:rem:u33e}
This answer is easy: the generalisations contain the configuration of Question~3.18 itself, which
was already answered negatively by Ackelsberg~\cite[Theorem~1.7]{ackelsberg2024counterexamples}
(the second arXiv version of~\cite{kra2023problems} does not record this), in exactly the way that
Questions~3.15 and~3.16 are answered by Remark~3.17 of~\cite{kra2023problems}. For (F2) the survey
itself already suffices, through the family $\{x^2\}$. We include it for completeness. A more
interesting question is which single polynomials $p\in\Z[x,y]$ have the property that every set of
positive upper Banach density contains a translate of $\{p(b_i,b_j)+b_k:i<j<k\}$ for some infinite $B$;
we do not know the answer, for instance, for $p(x,y)=x(y-x)$.
\end{remark}


\section{Sumsets of the form \texorpdfstring{$B+mB$}{B+mB}}\label{IA:sec:BmB}

This section concerns Subsection~3.4 of~\cite{kra2023problems}. For $B\subset\N$ and $m\in\N$ the
survey writes $B+mB=\{b+mb':b,b'\in B\}$; this unrestricted sumset contains the diagonal elements
$(m+1)b$. We write
\[
B\oplus mB=\{b+mb':\ b,b'\in B,\ b\ne b'\}
\]
for the restricted version. We answer Question~3.20 negatively, for two and for three or more
colours, even after excluding $m=1$ (\cref{IA:sec:q320}); Question~3.23 negatively as stated, the
answer depending on the diagonal (\cref{IA:sec:q323}); and Question~3.24 negatively
(\cref{IA:sec:q324}), where a counterexample is in fact already implicit in the work of Kousek and
Radi\'c~\cite{kousek2024infinite}.

\subsection{Colourings and Question~3.20}\label{IA:sec:q320}
A finite colouring of $\N$ is a map $\chi\colon\N\to\{0,1,\dots,r-1\}$. The survey gives, for each
$m$, a $3$-colouring with no monochromatic $B+mB$, $B$ infinite~\cite[Example~3.9]{kra2023problems}.
Two colours already suffice: Hindman~\cite[Theorem~2.11]{hindman1979partitions} proved that for every
$m\ge2$ the $2$-colouring $n\mapsto\lfloor\log_mn\rfloor\bmod2$ admits no infinite $B$ with
$\{mx+y:x,y\in B,\,x\ne y\}$ monochromatic (see also~\cite[Remark~1.5]{huang2026affirmative}).
Observing that such colourings depend heavily on $m$, the survey asks:

\begin{question}{3.20}\label{IA:q:320}
Is it true that for any finite colouring of $\N$ there exists $q\in\N$ such that for all $m\in\N$
with $\gcd(m,q)=1$ there is an infinite set $B\subset\N$ such that $B+mB$ is monochromatic?
\end{question}

The motivation given is that ``the set of problematic $m$'' might be ``small in a multiplicative
sense''. Since $\gcd(1,q)=1$ for every $q$, a positive answer for a colouring $\chi$ would require an
infinite $B$ with $B+B$ monochromatic. Hindman~\cite[Theorem~2.4]{hindman1979partitions} constructed
a $3$-colouring of $\N$ with no such $B$, so \cref{IA:q:320} has a negative answer as stated, for
three or more colours. This hinges on the literal formulation: for two colours, $m=1$ is harmless by
the recent solution of Owings' problem~\cite{owings1974problem} by Huang, Lian, Shao, Xiao, Xu and
Zhang~\cite[Theorem~1.1]{huang2026affirmative}, and the heuristic behind the question concerns
$m\ge2$. We show that the answer is negative also for $m\ge2$ and for two colours (for two colours,
with the unrestricted $B+mB$ of the survey). What is new is that a \emph{single} colouring is bad for
infinitely many $m$ simultaneously, including $m$ in every reduced residue class and infinitely many
primes; for each individual $m$, Hindman's theorem already provides a $2$-colouring for which $m$ is
strongly bad.

For a colouring $\chi$ call $m\in\N$ \emph{bad} if no infinite $B\subset\N$ has $B+mB$
monochromatic, and \emph{strongly bad} if no infinite $B$ has $B\oplus mB$ monochromatic.
Strongly bad integers are bad.

\begin{theorem}\label{IA:thm:320A}
Let $r\ge3$ and $\chi_r(n)=\lfloor\log_2n\rfloor\bmod r$. Every integer $m$ with $2^k<m<2^{k+1}$ for
some $k\ge1$ with $k\not\equiv0,-1\pmod r$ is strongly bad for $\chi_r$. The set of these $m$ has
logarithmic density $1-2/r$ and contains infinitely many primes.
\end{theorem}

\begin{theorem}\label{IA:thm:320B}
For $\rho>1$ let $\chi^{(\rho)}(n)=\lfloor\log_\rho n\rfloor\bmod 2$.
\begin{enumerate}[label=\textup{(\roman*)}]
\item\label{IA:it:lucas} Let $\varphi=\frac{1+\sqrt5}2$ and let $L_1=1$, $L_2=3$,
$L_{n+1}=L_n+L_{n-1}$ be the Lucas numbers. For every odd $n\ge3$, $L_n$ is bad for
$\chi^{(\varphi)}$. For every $q\in\N$ there are infinitely many odd $n\ge3$ with
$L_n\equiv1\pmod q$.
\item\label{IA:it:primes} There are $\rho\in(\frac32,2)$ and infinitely many primes that are bad for
$\chi^{(\rho)}$.
\end{enumerate}
\end{theorem}

\begin{corollary}\label{IA:cor:320}
For every $r\ge2$ there is an $r$-colouring of $\N$ such that for every $q\in\N$ some prime $m\ge2$
with $\gcd(m,q)=1$ is bad. In particular \cref{IA:q:320} has a negative answer for every number
$r\ge2$ of colours, also when the condition is imposed only for $m\ge2$ or only for prime $m$.
\end{corollary}

\begin{proof}
For $r\ge3$ use \cref{IA:thm:320A}: infinitely many primes are bad, and $q$ has finitely many prime
factors. For $r=2$ use \cref{IA:thm:320B}\ref{IA:it:primes} (or, for not necessarily prime $m$,
\cref{IA:thm:320B}\ref{IA:it:lucas}, where $L_n\equiv1\pmod q$ gives $\gcd(L_n,q)=1$).
\end{proof}

By \cref{IA:thm:320A}, the set of bad $m$ need not be small in any multiplicative sense: for $r$
colours it can have logarithmic density $1-2/r$. Question~3.19 of~\cite{kra2023problems}, whether
for every finite colouring some $m$ is not bad, remains open for $r\ge3$ colours; for $r=2$ it is
answered positively with $m=1$ by~\cite[Theorem~1.1]{huang2026affirmative}.

\begin{remark}[The diagonal and two colours]\label{IA:rem:320diag}
For $m\ge2$, $B+mB$ is the union of the diagonal $\{(m+1)b:b\in B\}$ and the two ordered pieces
$\{mx+y:x,y\in B,\,x<y\}$ and $\{mx+y:x,y\in B,\,x>y\}$. Huang et
al.~\cite[Theorem~1.6]{huang2026affirmative} proved that for every $2$-colouring of $\N$ and all
$m,\ell\in\N$ there is an infinite $B$ with $(m+\ell)B\cup\{mx+\ell y:x,y\in B,\,x<y\}$ monochromatic.
Applied with $(m,1)$ and with $(1,m)$, this shows that for two colours the diagonal together with
either ordered piece can always be made monochromatic. So a $2$-colouring for which some $m$ is bad
must use both ordered pieces; the proof of \cref{IA:thm:320B} uses one element from each of the three
parts. Hindman's theorem~\cite[Theorem~2.11]{hindman1979partitions} shows that each single $m\ge2$
is strongly bad for some $2$-colouring, namely $\lfloor\log_mn\rfloor\bmod2$ (this is the colouring
$\chi^{(\rho)}$ of \cref{IA:thm:320B} with $\rho=m$). Whether a single $2$-colouring can make
infinitely many pairwise coprime $m$ strongly bad, that is, whether Hindman's theorem has a version
uniform in $m$ (equivalently, the version of \cref{IA:q:320} with $B\oplus mB$ and two colours),
remains open. For related $2$- and $3$-colour obstructions to weighted ordered sums see
also~\cite[Propositions~5.1 and~5.2]{huang2026affirmative}.
\end{remark}

\begin{remark}
Di Nasso, Luperi Baglini, Mamino, Mennuni and Ragosta~\cite[\S1]{dinasso2025ramsey} answer negatively
the shift-free versions of several questions of Subsections~3.3 and~3.4 of~\cite{kra2023problems}
(in the numbering of its published version) in a strengthened ``Ramsey partition regular'' form.
These are different statements from \cref{IA:q:320}.
\end{remark}

For $\rho>1$ and $n\in\N$ put $\lambda_\rho(n)=\lfloor\log_\rho n\rfloor$, so that $\lambda_\rho(n)=k$
if and only if $\rho^k\le n<\rho^{k+1}$.

\begin{lemma}\label{IA:lem:floor}
If $x,y\in\R$ and $y-x\in(k,k+1)$ for some $k\in\Z$, then $\lfloor y\rfloor-\lfloor x\rfloor\in\{k,k+1\}$.
\end{lemma}

\begin{proof}
$\lfloor y\rfloor-\lfloor x\rfloor$ is an integer in $(y-x-1,y-x+1)\subset(k-1,k+2)$.
\end{proof}

\begin{lemma}\label{IA:lem:asym}
Let $\rho>1$, $m\ge2$ and $b,b'\in\N$. Put $\delta=\log_\rho(1+mb/b')$ and
$\delta'=\log_\rho(1+b/(mb'))$. Then $\log_\rho(b'+mb)=\log_\rho b'+\delta$,
$\log_\rho(mb'+b)=\log_\rho b'+\log_\rho m+\delta'$, and $0<\delta'<\delta$; moreover $\delta\to0$
as $b'\to\infty$ with $b$ fixed.
\end{lemma}

\begin{proof}
The identities are immediate, and $b/(mb')<mb/b'$ because $m\ge2$.
\end{proof}

\begin{proposition}\label{IA:prop:320A}
Let $\rho>1$, $r\ge3$ and $\chi(n)=\lambda_\rho(n)\bmod r$. Let $m\ge2$ be such that
$L=\log_\rho m\notin\Z$ and $k=\lfloor L\rfloor$ satisfies $k\not\equiv0$ and $k+1\not\equiv0
\pmod r$. Then $m$ is strongly bad for $\chi$.
\end{proposition}

\begin{proof}
Let $B$ be infinite and fix $b\in B$. By \cref{IA:lem:asym} there is $b'\in B$, $b'\ne b$, with
$\delta<L-k$. Put $X=b'+mb$ and $Y=b+mb'$, both in $B\oplus mB$. By \cref{IA:lem:asym},
$\log_\rho Y-\log_\rho X=L+\delta'-\delta\in(L-\delta,L)\subset(k,k+1)$, so by \cref{IA:lem:floor}
$\lambda_\rho(Y)-\lambda_\rho(X)\in\{k,k+1\}$. Neither number is divisible by $r$, so
$\chi(X)\ne\chi(Y)$.
\end{proof}

\begin{proof}[Proof of \cref{IA:thm:320A}]
For $2^k<m<2^{k+1}$ we have $\log_2m\notin\Z$ and $\lfloor\log_2m\rfloor=k$, so
\cref{IA:prop:320A} with $\rho=2$ applies. Let $S$ be the set of these $m$. For each $k$,
$\sum_{2^k<m<2^{k+1}}1/m=\log2+O(2^{-k})$, and among any $r$ consecutive values of $k$ exactly $r-2$
are admissible. Hence $\sum_{m\in S,\,m<2^{K+1}}1/m=\frac{r-2}r K\log2+O(1)$, and for $2^K\le
N<2^{K+1}$ the partial sum up to $N$ differs from this by $O(1)$. So $S$ has logarithmic density
$1-2/r$. By Bertrand's postulate each interval $(2^k,2^{k+1})$, $k\ge1$, contains a prime.
\end{proof}

\begin{proposition}\label{IA:prop:320B}
Let $m\ge2$ and $\rho>1+\frac1m$, and suppose there is an odd integer $j$ with $m<\rho^j<m+1$. Then
$m$ is bad for $\chi(n)=\lambda_\rho(n)\bmod2$.
\end{proposition}

\begin{proof}
Let $L=\log_\rho m$ and $L'=\log_\rho(m+1)$. Then $L<j<L'$ and $L'-L=\log_\rho(1+\frac1m)<1$. Let
$\eta=\min(j-L,L'-j)>0$. Let $B$ be infinite, fix $b\in B$, and by \cref{IA:lem:asym} choose $b'\in
B$, $b'\ne b$, with $\delta<\eta$. Consider the elements
\[
X=b'+mb,\qquad Y=b+mb',\qquad Z=b'+mb'=(m+1)b'
\]
of $B+mB$, and let $u=\log_\rho b'$, $x=\log_\rho X$, $y=\log_\rho Y$, $z=\log_\rho Z$. By
\cref{IA:lem:asym}, $x=u+\delta$, $y=u+L+\delta'$ and $z=u+L'$ with $0<\delta'<\delta<\eta$. Hence
$y<u+L+\eta\le u+j<x+j$ and $z=u+L'\ge u+j+\eta>x+j$, so
\begin{equation}\label{IA:eq:sandwich}
y<x+j<z,\qquad 0<z-y<L'-L<1 .
\end{equation}
Suppose $X,Y,Z$ have the same colour. Since $0<z-y<1$, the integer $\lambda_\rho(Z)-\lambda_\rho(Y)$ is
$0$ or $1$, and it is even; so $\lambda_\rho(Y)=\lambda_\rho(Z)=:N$. Then $N\le y<z<N+1$, and
by~\eqref{IA:eq:sandwich} $N<x+j<N+1$, so $N=\lfloor x\rfloor+j=\lambda_\rho(X)+j$. As $j$ is odd,
$\chi(X)\ne\chi(Y)$, a contradiction.
\end{proof}

\begin{lemma}\label{IA:lem:nested}
There are $\rho\in(\frac32,2)$, odd integers $j_1<j_2<\cdots$ and primes $p_1<p_2<\cdots$ with
$p_k<\rho^{j_k}<p_k+1$ for all $k$.
\end{lemma}

\begin{proof}
We construct closed intervals $[\frac32,2]=I_0\supset I_1\supset\cdots$, odd $j_k$ and primes $p_k$
with $p_k<s^{j_k}<p_k+1$ for all $s\in I_k$; any $\rho\in\bigcap_kI_k$ then works. Let $p_0=1$,
$j_0=0$, and suppose $I_{k-1}=[\alpha,\beta]$ with $\frac32\le\alpha<\beta$ is constructed. Choose an
odd $j>j_{k-1}$ so large that $x:=\alpha^j\ge p_{k-1}+2$ and $\beta^j\ge4x$, put $n=\lceil x\rceil$,
and let $p$ be a prime with $n<p<2n$ (Bertrand's postulate). Then $\alpha^j\le n<p$, $p>p_{k-1}$,
and $p+1\le2n\le2x+2<4x\le\beta^j$. Let $I_k=[(p+\frac13)^{1/j},(p+\frac23)^{1/j}]\subset(\alpha,\beta)$,
$j_k=j$ and $p_k=p$.
\end{proof}

\begin{proof}[Proof of \cref{IA:thm:320B}]
(i) Let $\psi=\frac{1-\sqrt5}2=-\varphi^{-1}$, so $L_n=\varphi^n+\psi^n$. For odd $n$,
$\psi^n\in(-1,0)$, so $L_n<\varphi^n<L_n+1$. For $n\ge3$, $L_n\ge4$ and $\varphi>1+\frac1{L_n}$, so
\cref{IA:prop:320B} applies with $m=L_n$, $\rho=\varphi$, $j=n$. For the second claim, $(L_n\bmod q)$
satisfies an invertible linear recurrence over $\Z/q\Z$ and is therefore purely periodic, with some
period $P$; then $L_n\equiv L_1=1\pmod q$ for all $n\equiv1\pmod{2P}$, and these $n$ are odd.

(ii) Take $\rho$, $(j_k)$, $(p_k)$ from \cref{IA:lem:nested}. Then $\rho>\frac32\ge1+\frac1{p_k}$ and
$p_k<\rho^{j_k}<p_k+1$ with $j_k$ odd, so \cref{IA:prop:320B} applies to $m=p_k$.
\end{proof}

\formalization{\Cref{IA:thm:320A,IA:thm:320B}, \cref{IA:cor:320} and the propositions and
lemmas they rely on are formalized in \leanfile{Q320} (\leanname{thm320A}, \leanname{thm320B},
\leanname{cor320}). Hindman's $3$-colouring, quoted for $m=1$, is not formalized.}


\subsection{Logarithmic density and Question~3.23}\label{IA:sec:q323}
The survey shows that a density version of Question~3.19 needs care: there is a thick set of positive
upper density containing no $B\oplus mB+t$ with $m\ge2$ and $B$ infinite
\cite[Example~3.21]{kra2023problems}. That set has zero upper logarithmic density, which motivates:

\begin{question}{3.23}\label{IA:q:323}
Does every set $A\subset\N$ with $\udlog(A)>0$ contain a sumset $B+mB+t$ for some $m\ge2$, $t\ge0$
and some infinite set $B\subset\N$?
\end{question}

As written, $B+mB$ is the unrestricted sumset: the survey writes $B\oplus B$ or states $b_1\ne b_2$
explicitly when it means the restricted one, and calls $B+B$ the unrestricted sumset. The answer to
\cref{IA:q:323} as stated is negative.

For each \emph{fixed} pair $(\ell,m)$, unrestricted sums $mB+\ell B$ have been studied by
Kousek~\cite{kousek2025asymmetric}. He showed that $\ldens(A)>1/2$ forces
$\{mb_1+\ell b_2:b_1\le b_2\}+t\subset A$, and that a suitable upper density condition depending on
$m/\ell$ does too~\cite[Theorems~1.6 and~1.7]{kousek2025asymmetric}; these are sharp, and for distinct
$\ell,m$ there are sets of positive lower density containing no shift of
$mB+\ell B=\{mb_1+\ell b_2:b_1,b_2\in B\}$ with $B$ infinite~\cite[Propositions~5.1, 6.5, 7.6
and~7.7]{kousek2025asymmetric}; see also his Questions~7.8 and~7.9. Those examples depend on
$(\ell,m)$, whereas \cref{IA:q:323} allows $m$ to depend on $A$. What is new below is a single set of
positive logarithmic density that avoids $\ell B+mB+t$ for all $\ell,m\ge1$ and $t\in\Z$ at once. As in
Kousek's results (and as in \cref{IA:rem:320diag} for colourings), the diagonal together with one
ordered piece cannot be avoided by sets of lower density $>1/2$, and our proof uses the diagonal and
both ordered pieces.

Fix $\eta\in(0,\frac16)$ and let $U_\eta=(\frac13+\eta,\frac23-\eta)$. Let $T_k=2^{2^k}$,
$X_k=kT_k$ and $Y_k=T_{k+1}/k=T_k^2/k$ for $k\ge1$; note $T_k\le X_k<Y_k\le T_{k+1}$. Let
\[
R_k=\{s\in\{0,\dots,k-1\}:\ s/k\in U_\eta\},\qquad
A_\eta=\bigcup_{k\ge1}\{n\in\N:\ X_k\le n<Y_k,\ n\bmod k\in R_k\}.
\]
Thus on the block $[X_k,Y_k)$ the set $A_\eta$ is the periodic set $\{n:\{n/k\}\in U_\eta\}$, and
the transition intervals $[Y_{k-1},X_k)$ are empty.

\begin{theorem}\label{IA:thm:323}
\begin{enumerate}[label=\textup{(\alph*)}]
\item\label{IA:it:323dens} The logarithmic density of $A_\eta$ exists and equals $\frac13-2\eta$.
\item\label{IA:it:323avoid} For all integers $\ell,m\ge1$, all $t\in\Z$ and every infinite
$B\subset\N$, we have $\ell B+mB+t\not\subset A_\eta$. More precisely, for every $c\in B$ only
finitely many $b\in B$ satisfy
\[
\ell b+mc+t,\quad \ell c+mb+t,\quad (\ell+m)b+t\ \in A_\eta .
\]
\end{enumerate}
In particular (with $\ell=1$) $A_\eta$ is a negative answer to \cref{IA:q:323}, simultaneously for
all $m$.
\end{theorem}

\begin{remark}\label{IA:rem:323}
\begin{enumerate}[label=(\roman*)]
\item The answer hinges on the diagonal: the contradiction comes from the element $(\ell+m)b+t$.
The example motivating \cref{IA:q:323}, \cite[Example~3.21]{kra2023problems}, concerns the
restricted sumset, and for it $A_\eta$ is no counterexample: for every odd $m\ge3$ there is an
infinite $B$ with $B\oplus mB\subset A_\eta$. Indeed, let $\eps=\frac12(\frac16-\eta)$, so that
$[\frac12,\frac12+\eps)\subset U_\eta$. Choose $b_1<b_2<\cdots$ inductively: given $b_1,\dots,b_{j-1}$,
choose an even $k$ so large that $(m+1)b_{j-1}<\eps k$ and $Y_k/(m+1)-X_k>k$, and choose
$b_j\in[X_k,Y_k/(m+1))$ with $b_j\equiv k/2\pmod k$. For $i<j$ the numbers $b_j+mb_i$ and
$mb_j+b_i$ lie in $[X_k,Y_k)$ and are congruent modulo $k$ to $\frac k2+mb_i$ and to
$\frac k2+b_i$ respectively (as $m$ is odd, $\frac{mk}2\equiv\frac k2$), both of which lie in
$[\frac k2,\frac k2+\eps k)$. Hence both belong to $A_\eta$.
\item The restricted version of \cref{IA:q:323} (does every $A$ with $\udlog(A)>0$ contain
$B\oplus mB+t$ for some $m\ge2$, $t\ge0$ and infinite $B$?) remains open.
\end{enumerate}
\end{remark}

\begin{lemma}\label{IA:lem:sumfree}
Let $k\ge1$, $s_1,s_2,s_3\in R_k$ and $s\in\Z$ with $|s|<\eta k$. Then $s_1+s_2-s_3\not\equiv s\pmod k$.
\end{lemma}

\begin{proof}
From $s_i/k\in U_\eta$ we get $s_1+s_2-s_3-s\in\big(k(\frac23+2\eta)-k(\frac23-\eta)-\eta k,\
k(\frac43-2\eta)-k(\frac13+\eta)+\eta k\big)=(2\eta k,\,k(1-2\eta))\subset(0,k)$, which contains no
multiple of $k$.
\end{proof}

\begin{lemma}\label{IA:lem:blocksep}
If $n,n'\in A_\eta$ lie in the blocks of indices $k<k'$, then $n'/n>k(k+1)$. If $n\in A_\eta$ and
$n\ge T_{K+1}$, then the block of $n$ has index at least $K+1$.
\end{lemma}

\begin{proof}
$n<Y_k=T_{k+1}/k$ and $n'\ge X_{k+1}=(k+1)T_{k+1}$. The block of index $k$ lies in $[T_k,T_{k+1})$.
\end{proof}

\begin{lemma}\label{IA:lem:perblock}
Let $k\in\N$, $R\subset\{0,\dots,k-1\}$, $\rho=|R|/k$, and $1\le P<Q$ real. Let
$g\colon\Z\cap[P,Q)\to\{0,1\}$ be constant on each of $J$ consecutive subintervals partitioning
$[P,Q)$. Then
\[
\Big|\sum_{\substack{P\le n<Q\\ n\bmod k\in R}}\frac{g(n)}n-\rho\sum_{P\le n<Q}\frac{g(n)}n\Big|\le\frac{2Jk}P .
\]
\end{lemma}

\begin{proof}
It suffices to show that the error on one subinterval $[P_1,Q_1)$, $P_1\ge P$, on which $g\equiv1$,
is at most $2k/P$. Let $n_0=\lceil P_1\rceil$ and split the integers of $[P_1,Q_1)$ into complete
periods $\{n_0+ik,\dots,n_0+(i+1)k-1\}$, $0\le i<I$, and a remainder of fewer than $k$ integers. On
the $i$-th period both sums lie in $[\rho k/(n_0+(i+1)k),\,\rho k/(n_0+ik)]$, so the errors telescope
to at most $\rho k/n_0\le k/P$. On the remainder both sums lie in $[0,k/n_0]$.
\end{proof}

\begin{proof}[Proof of \cref{IA:thm:323}\ref{IA:it:323avoid}]
Fix $\ell,m\ge1$, $t\in\Z$ and $c\in B$, and put $s=(\ell+m)c+t$ and $C=\ell+m+1$. For $b\in B$ let
\[
u=\ell b+mc+t,\qquad v=\ell c+mb+t,\qquad w=(\ell+m)b+t,
\]
all three elements of $\ell B+mB+t$, and note the identity $u+v-w=s$. If
$b\ge2(|t|+(\ell+m)c)$, then $u,v,w\in[b/2,Cb]$, so their pairwise ratios are at most $2C$. Choose
$K$ with $K(K+1)>2C$ and $K\ge|s|/\eta$, and suppose $b\ge\max\big(2(|t|+(\ell+m)c),\,2T_{K+1}\big)$ and
$u,v,w\in A_\eta$. By \cref{IA:lem:blocksep} the blocks of $u,v,w$ have indices $\ge K+1$, and they
coincide, since otherwise two of the numbers would have ratio $>(K+1)(K+2)>2C$. Let $k>K$ be the common
index. The residues of $u,v,w$ modulo $k$ lie in $R_k$ and satisfy $(u\bmod k)+(v\bmod k)-(w\bmod k)
\equiv s\pmod k$ with $|s|<\eta k$, contradicting \cref{IA:lem:sumfree}. So only finitely many $b\in B$
have $u,v,w\in A_\eta$, and as $B$ is infinite, $\ell B+mB+t\not\subset A_\eta$.
\end{proof}

\begin{proof}[Proof of \cref{IA:thm:323}\ref{IA:it:323dens}]
All logarithms are natural. Let $H(P,Q)=\sum_{P\le n<Q}1/n$, so that
$H(P,Q)=\log(Q/P)+O(1/P)$ for $1\le P<Q$. Let $\rho_k=|R_k|/k$ and $\mu=|U_\eta|=\frac13-2\eta$. The open interval
$(k(\frac13+\eta),k(\frac23-\eta))$ has length $k\mu$, so $|\rho_k-\mu|\le1/k$.

Let $T_K\le N<T_{K+1}$ and $S(N)=\sum_{n\le N,\,n\in A_\eta}1/n$. By \cref{IA:lem:perblock} with
$J=1$,
\[
S(N)=\sum_{k<K}\rho_kH(X_k,Y_k)+\rho_K H\big(X_K,\min(N+1,Y_K)\big)+O\Big(\sum_k\frac{2k}{X_k}\Big),
\]
where the middle term is read as $0$ if $N<X_K$, and $\sum_k2k/X_k=\sum_k2/T_k<\infty$. We have
$H(X_k,Y_k)=2^k\log2-2\log k+O(1/T_k)$, and
$\log N=2\log 2+\sum_{1\le k<K}2^k\log2+\log(N/T_K)$. Moreover the middle term differs from
$\mu\log(N/T_K)$ by at most $|\rho_K-\mu|2^K\log 2+O(\log K)$: if $N<X_K$ then $\log(N/T_K)<\log K$;
if $X_K\le N<Y_K$ then $H(X_K,N+1)=\log(N/T_K)-\log K+O(1)$; and if $N\ge Y_K$ then both
$H(X_K,Y_K)$ and $\log(N/T_K)$ lie within $2\log K+O(1)$ of $2^K\log 2$. Altogether
\[
\big|S(N)-\mu\log N\big|\le\sum_{k\le K}\frac{2^k\log2}{k}+O(K\log K)=O(2^K/K),
\]
while $\log N\ge2^K\log2$. Hence $S(N)/\log N\to\mu$.
\end{proof}

\formalization{Both parts of \cref{IA:thm:323} are formalized in \leanfile{Q323}
(\leanname{thm323}); \cref{IA:rem:323} is not.}


\subsection{Logarithmic equidistribution and Question~3.24}\label{IA:sec:q324}
All examples of the survey of sets of positive density avoiding some unrestricted sumset $B+B+t$
are built at a logarithmic scale; to test whether this is the only obstruction, the survey asks:

\begin{question}{3.24}\label{IA:q:324}
If $A\subset\N$ has $\dlog(A)>0$ and for every interval $(a,b)\subset[0,1)$ and every $\theta>0$
\begin{equation}\label{IA:eq:324}
\dlog\big(\{n\in A:\ \{\theta\log n\}\in(a,b)\}\big)=\dlog(A)\cdot(b-a),
\end{equation}
does $A$ contain a sumset $B+B+t$ for some $t\ge0$ and some infinite set $B\subset\N$?
\end{question}

The answer is negative, and a counterexample is already implicit in the work of Kousek and
Radi\'c~\cite{kousek2024infinite}: the set from the first part of~\cite[Proposition~4.2]{kousek2024infinite},
which has natural density $\frac12$ and contains no $B+B+t$, also satisfies~\eqref{IA:eq:324}, as we
check in \cref{IA:prop:KR} (this observation is not made in~\cite{kousek2024infinite}). By
\cite[Theorem~1.3]{kousek2024infinite}, every set of lower density $>\frac12$ contains some $B+B+t$,
so the value $\frac12$ is optimal among sets having a natural density; we do not know whether
$\dlog(A)>\frac12$ together with~\eqref{IA:eq:324} forces $B+B+t\subset A$. That example couples a
logarithmic phase with residue classes modulo $2$. In \cref{IA:thm:324} we give an alternative
example with no residue-class coupling at all, in which a single logarithmic phase drifts slowly;
it shows that the test~\eqref{IA:eq:324}, which involves only linear phases $\theta\log n$, does
not detect logarithmic-scale obstructions. Recall that by Kra, Moreira, Richter
and Robertson~\cite[Theorem~1.2]{kra2022proof} every set of positive upper Banach density contains
$B\oplus B+t$ for some infinite $B$ and $t\in\N$; so all these examples necessarily exploit the
diagonal elements $2b$.

\begin{lemma}\label{IA:lem:logsum}
Let $F\subset[0,\infty)$ be a union of pairwise disjoint intervals with left endpoints $\alpha_i$
such that $\#\{i:\alpha_i\le L\}\le C(L+1)$ for all $L\ge0$. Let $k\in\N$ and $j\in\Z$. Then, as
$N\to\infty$,
\[
\sum_{\substack{n\le N\\ \log_2n\in F}}\frac1n=\log2\cdot\big|F\cap[0,\log_2N]\big|+O(1)
\qquad\text{and}\qquad
\sum_{\substack{n\le N,\ n\equiv j\ (k)\\ \log_2n\in F}}\frac1n
=\frac1k\sum_{\substack{n\le N\\ \log_2n\in F}}\frac1n+O(1),
\]
where $|\cdot|$ is Lebesgue measure.
\end{lemma}

\begin{proof}
If $1\le x\le y$ and $I$ is an interval with endpoints $x,y$, then $|\sum_{n\in I}1/n-\log(y/x)|\le
1/x$, by comparison with $\int dt/t$. Each component of $F\cap[0,\log_2N]$ with left endpoint
$\alpha_i$ corresponds to such an interval $I$ with $x=2^{\alpha_i}$, and $\log(y/x)=\log2$ times
the length of the component. The total error is at most $\sum_i2^{-\alpha_i}\le\sum_{L\ge0}
C(L+1)2^{1-L}<\infty$. The second claim follows in the same way from \cref{IA:lem:perblock} ($J=1$),
applied to each component, with total error at most $\sum_i 2k\,2^{-\alpha_i}<\infty$.
\end{proof}

\begin{proposition}[set of Kousek--Radi\'c]\label{IA:prop:KR}
Let $A^{(1)}=\bigcup_{n\ge1}[4^n,(2-\frac1n)4^n)$, $A^{(2)}=\bigcup_{n\ge1}[(2+\frac1n)4^n,4^{n+1})$
and $A_{\mathrm{KR}}=(A^{(1)}\cap2\N)\cup(A^{(2)}\cap(2\N+1))$ (intersected with $\N$). Then
$A_{\mathrm{KR}}$ has natural density $\frac12$, contains no $B+B+t$ with $t\ge0$ and $B$ infinite,
and satisfies~\eqref{IA:eq:324} with $\dlog(A_{\mathrm{KR}})=\frac12$.
\end{proposition}

\begin{proof}
The first two assertions are~\cite[Proposition~4.2]{kousek2024infinite}, stated there for $t\in\N$;
the case $t=0$ follows, because if $B+B\subset A_{\mathrm{KR}}$ then $B'=\{b-1:b\in B,\ b\ge2\}$
satisfies $B'+B'+2\subset A_{\mathrm{KR}}$.

For~\eqref{IA:eq:324}, let $V_0=[0,\frac12)$, $V_1=[\frac12,1)$ and
$A^*=\{n\in2\N:\{\frac12\log_2n\}\in V_0\}\cup\{n\in2\N+1:\{\frac12\log_2n\}\in V_1\}$. Since
$\{\frac12\log_2n\}\in V_0$ if and only if $n\in[4^j,2\cdot4^j)$ for some $j\ge0$, the symmetric
difference of $A_{\mathrm{KR}}$ and $A^*$ is contained in $[1,4)\cup W$, where
$W=\bigcup_{n\ge1}[(2-\frac1n)4^n,(2+\frac1n)4^n]$. For $N<4^{K+1}$ we have
$\sum_{n\le N,\,n\in W}\frac1n\le\sum_{n\le K}\big(\log\frac{2+1/n}{2-1/n}+4^{-n}\big)=O(\log K)
=o(\log N)$, so it suffices to verify~\eqref{IA:eq:324} for $A^*$. Fix $\theta>0$ and
$(a,b)\subset[0,1)$, put $\theta'=\theta\log2$, so that $\{\theta\log n\}=\{\theta'\log_2n\}$, and
let $F_i=\{x\ge0:\{x/2\}\in V_i,\ \{\theta'x\}\in(a,b)\}$ for $i\in\{0,1\}$. These sets satisfy the
hypothesis of \cref{IA:lem:logsum}, and $F_0\sqcup F_1=F:=\{x\ge0:\{\theta'x\}\in(a,b)\}$. By
\cref{IA:lem:logsum},
\[
\sum_{\substack{n\le N,\ n\in A^*\\ \{\theta\log n\}\in(a,b)}}\frac1n
=\sum_{i=0}^1\ \sum_{\substack{n\le N,\ n\equiv i\ (2)\\ \log_2n\in F_i}}\frac1n
=\frac12\sum_{i=0}^1\ \sum_{\substack{n\le N\\ \log_2n\in F_i}}\frac1n+O(1)
=\frac12\log 2\,|F\cap[0,\log_2N]|+O(1),
\]
and $|F\cap[0,L]|=(b-a)L+O(1)$. Dividing by $\log N=\log2\cdot\log_2N$ gives the limit
$\frac12(b-a)$. The same computation without the condition on $\{\theta\log n\}$ (that is, with
$F=[0,\infty)$) gives $\dlog(A^*)=\frac12$.
\end{proof}

The colour class $(A_0\cap3\N)\cup(A_1\cap(3\N+1))\cup(A_2\cap(3\N+2))$, with
$A_i=\{n:\{\frac12\log_2n\}\in[\frac i3,\frac{i+1}3)\}$, from the remark after
Proposition~5.1 of~\cite{kousek2024infinite} is handled in the same way (with residues modulo $3$)
and gives a counterexample of logarithmic density $\frac13$.

\begin{theorem}\label{IA:thm:324}
Let $0<\eta<\frac12$, $f(x)=\frac x2+\sqrt x$ for $x\ge0$, and
\[
A=\big\{n\in\N:\ \{f(\log_2n)\}\in[0,\tfrac12-\eta)\big\}.
\]
Then $\dlog(A)=\frac12-\eta$, $A$ satisfies~\eqref{IA:eq:324}, and $A$ contains no $B+B+t$ with
$t\in\Z$ and $B\subset\N$ infinite.
\end{theorem}

\begin{lemma}\label{IA:lem:oscint}
Let $\theta'>0$ and $(h,k)\in\Z^2\setminus\{(0,0)\}$. Then
$\frac1L\int_0^L\e\big(hf(x)+k\theta'x\big)\,dx\to0$ as $L\to\infty$, where $\e(y)=e^{2\pi iy}$.
\end{lemma}

\begin{proof}
The phase is $\psi(x)=\gamma x+h\sqrt x$ with $\gamma=\frac h2+k\theta'$. If $\gamma\ne0$, let
$x_0=(h/\gamma)^2$. On $[x_0,L]$ the derivative $\psi'(x)=\gamma+\frac h{2\sqrt x}$ is monotone and
$|\psi'|\ge|\gamma|/2$. Integrating by parts,
\[
\int_{x_0}^L\e(\psi)\,dx=\Big[\frac{\e(\psi)}{2\pi i\psi'}\Big]_{x_0}^L
-\int_{x_0}^L\e(\psi)\,d\Big(\frac1{2\pi i\psi'}\Big),
\]
and since $1/\psi'$ is monotone with $|1/\psi'|\le2/|\gamma|$, each of the two terms is at most
$\frac{1}{2\pi}\cdot\frac{4}{|\gamma|}$.
The integral over $[0,x_0]$ is at most $x_0$. If $\gamma=0$, then $h\ne0$ (otherwise $k\theta'=0$, so
$k=0$), and $\int_0^L\e(h\sqrt x)\,dx=\int_0^{\sqrt L}2y\,\e(hy)\,dy=O(\sqrt L/|h|)$ by one
integration by parts. In both cases the integral is $o(L)$.
\end{proof}

\begin{proof}[Proof of \cref{IA:thm:324}]
Let $c=\frac12-\eta$, fix $\theta>0$, put $\theta'=\theta\log2$, and let
$\Psi(x)=(f(x)\bmod1,\ \theta'x\bmod1)\in\T^2$. By \cref{IA:lem:oscint} and linearity,
$\frac1L\int_0^LP(\Psi(x))\,dx\to\int_{\T^2}P$ for every trigonometric polynomial $P$, hence, by
uniform density of trigonometric polynomials in $C(\T^2)$, for every continuous $P$. If
$\Omega\subset\T^2$ is a box, then for each $\delta>0$ there are continuous $P_-\le\one_\Omega\le P_+$
with $\int(P_+-P_-)<\delta$, and therefore $\frac1L|\{x\in[0,L]:\Psi(x)\in\Omega\}|\to|\Omega|$. We
apply this to $\Omega=[0,c)\times(a,b)$ and to $\Omega=[0,c)\times\T$.

Since $f$ is continuous and strictly increasing with $f(0)=0$, the sets
$F=\{x\ge0:\{f(x)\}\in[0,c)\}=\bigcup_{j\ge0}[f^{-1}(j),f^{-1}(j+c))$ and
$F'=F\cap\{x:\{\theta'x\}\in(a,b)\}$ are unions of disjoint intervals; the endpoints of their
components in $[0,L]$ are among the at most $2(f(L)+1)+2(\theta'L+1)$ points where $f(x)$ or
$\theta'x$ hits $\Z\cup(\Z+c)$ or $(\Z+a)\cup(\Z+b)$, so \cref{IA:lem:logsum} applies. Since
$A=\{n:\log_2n\in F\}$ and $\{n\in A:\{\theta\log n\}\in(a,b)\}=\{n:\log_2n\in F'\}$, we obtain
$\dlog(A)=\lim|F\cap[0,L]|/L=c$ and the logarithmic density of $\{n\in A:\{\theta\log n\}\in(a,b)\}$
equals $\lim|F'\cap[0,L]|/L=c(b-a)$. This proves~\eqref{IA:eq:324}.

Now suppose $B$ is infinite, $t\in\Z$ and $B+B+t\subset A$. Let $b'=\min B$ and, for $b\in B$ large,
put $n_1=b+b'+t$, $n_2=2b+t$ (both in $B+B+t$) and $x_i=\log_2n_i$. As $b\to\infty$,
$x_2-x_1=\log_2\frac{2b+t}{b+b'+t}\to1$ and $|\sqrt{x_2}-\sqrt{x_1}|\le\frac{|x_2-x_1|}{2\sqrt{\min(x_1,x_2)}}
\to0$, so $\Delta_b:=f(x_2)-f(x_1)-\frac12\to0$. Take $b$ with $|\Delta_b|<\eta/2$. Since $n_1\in A$,
$\{f(x_1)\}\in[0,c)$, so $f(x_2)\equiv\{f(x_1)\}+\frac12+\Delta_b$ with
$\{f(x_1)\}+\frac12+\Delta_b\in[\frac12-\frac\eta2,1-\frac\eta2)$. Hence
$\{f(x_2)\}\notin[0,c)$, that is, $n_2\notin A$, a contradiction.
\end{proof}

\formalization{\Cref{IA:prop:KR}, including the two assertions quoted
from~\cite{kousek2024infinite}, and \cref{IA:thm:324} are formalized in \leanfile{Q324}
(\leanname{prop_KR}, \leanname{thm324}).}

\begin{remark}\label{IA:rem:324}
\begin{enumerate}[label=(\roman*)]
\item The obstruction in \cref{IA:thm:324} is still of logarithmic scale: the phase $f(\log_2n)$
increases by $\frac12+o(1)$ when $n$ is doubled. Equidistribution against all linear phases
$\theta\log n$ does not rule out such phases, because the term $\sqrt{\log_2n}$ makes $f$
asymptotically independent of every linear phase, including the resonant frequency
$\theta=1/(2\log2)$.
\item The set $A_\eta$ of \cref{IA:sec:q323} is a further counterexample to \cref{IA:q:324}: it
contains no $B+B+t$ by \cref{IA:thm:323}\ref{IA:it:323avoid} with $\ell=m=1$, and it
satisfies~\eqref{IA:eq:324}. To see the latter, let $g(n)=\one[\{\theta\log n\}\in(a,b)]$. On the
block $[X_k,Y_k)$, $g$ is constant on $J_k\le2\theta\log(Y_k/X_k)+3=O(2^k)$ subintervals, so by
\cref{IA:lem:perblock} the sum of $g(n)/n$ over $A_\eta\cap[X_k,Y_k)$ equals $\rho_k$ times the sum
over $[X_k,Y_k)$ up to $O(2^k/T_k)$, which is summable. The rest of the proof of
\cref{IA:thm:323}\ref{IA:it:323dens} applies verbatim with weights $g(n)$ (the transition intervals
contribute $O(K\log K)$ and the errors $|\rho_k-\mu|$ contribute $O(2^K/K)$), giving
$\sum_{n\le N,\,n\in A_\eta}g(n)/n=\mu\sum_{n\le N}g(n)/n+o(\log N)$; and
$\sum_{n\le N}g(n)/n=(b-a)\log N+O(1)$ by \cref{IA:lem:logsum}. Here the obstruction comes from scale-separated blocks with periodic structure inside each block,
rather than from a phase of $\log n$; the empty transition intervals $[Y_{k-1},X_k)$ have
logarithmic density zero and are invisible to~\eqref{IA:eq:324}.
\end{enumerate}
\end{remark}
